\section{Binary injectivity over an arbitrary Noetherian base}
\label{sec:binary-injective-classification}

This section proves that binary tensor injectivity determines every
polynomial tensor map over a Noetherian base. Let $R\subset S$ be a
unital inclusion of commutative rings, with
$R\ne0$.  Put
\[
 \Omega_n(R,S)=\{F\in S[X_1,\ldots,X_n]:F(R^n)\subseteq R\},
 \qquad \Omega(R,S)=\Omega_1(R,S),
\]
and let
\[
 \theta_n:\Omega(R,S)^{\otimes_R n}\longrightarrow\Omega_n(R,S)
\]
multiply one-variable polynomials in distinct variables.  All tensor
products in this section are algebraic tensor products.  Coefficients
in $S$ may have $R$-torsion.

\begin{theorem}[Binary injectivity and the local classification]
\label{thm:binary-injective-classification}
Suppose that $R$ is Noetherian.  For every maximal ideal
$\mathfrak m$ with finite residue field, define the subring
\[
 C_{\mathfrak m}
 =\bigcup_{j\ge0}
   (R_{\mathfrak m}:_{S_{\mathfrak m}}
                    (\mathfrak mR_{\mathfrak m})^j).
\]
The following conditions are equivalent.
\begin{enumerate}[label=\textup{(\roman*)}]
\item $\theta_2$ is injective.
\item $\theta_n$ is an isomorphism for every $n\ge1$.
\item For every $F\in\Omega(R,S)$, each of $F(X+Y)$ and $F(XY)$
      has exactly one preimage under $\theta_2$.
\item At each maximal ideal $\mathfrak m$ with finite residue field,
      either $C_{\mathfrak m}=R_{\mathfrak m}$, or
      $R_{\mathfrak m}$ is a one-dimensional domain whose completion
      is a domain and
      \[
       C_{\mathfrak m}=\operatorname{Frac}(R_{\mathfrak m})
       \subseteq S_{\mathfrak m}.
      \]
      The last inclusion is the canonical one extending
      $R_{\mathfrak m}\subseteq S_{\mathfrak m}$.
\end{enumerate}
There is no additional condition at maximal ideals with infinite
residue field.  When these conditions hold, $\Omega(R,S)$ has a unique
compatible $R$-plethory structure.
\end{theorem}

The implication \textup{(iii)}$\Rightarrow$\textup{(i)} already follows
from uniqueness for the zero polynomial.  The substantive necessity
statement is that binary injectivity alone forces \textup{(iv)}.
We first establish the base-change statements and the tensor
obstructions needed for that implication.

\Needspace{8\baselineskip}
\subsection{Base change and the effective coefficients}

\begin{lemma}[Localization and completion with arbitrary coefficients]
\label{lem:arbitrary-coefficient-basechange}
Let $R$ be Noetherian and $U\subset R$ multiplicatively closed.  Then
\[
 U^{-1}\Omega_n(R,S)
 =\Omega_n(U^{-1}R,U^{-1}S).
\]
If $(R,\mathfrak m)$ is local and $A=\widehat R$, then
\[
 A\otimes_R\Omega_n(R,S)
 =\Omega_n(A,A\otimes_R S).
\]
These identifications respect the maps $\theta_n$ and polynomial
substitution.  For a local base with finite residue field, every
coefficient occurring in any $\Omega_n(R,S)$ lies in
\[
 C=\bigcup_{j\ge0}(R:_S\mathfrak m^j).
\]
Moreover,
\[
 A\otimes_R C
 =\bigcup_{j\ge0}
   (A:_{A\otimes_R S}(\mathfrak mA)^j).
\]
For a local base with infinite residue field,
$\Omega_n(R,S)=R[X_1,\ldots,X_n]$.
\end{lemma}
\begin{proof}
For the forward localization inclusion, induct on total degree.
If $F$ has total degree $d$ and $s\in U$, then
$s^dF(\mathbf X)-F(s\mathbf X)$ has smaller total degree and still
takes $R^n$ into $R$.  Evaluating at $\mathbf a/s$, with
$\mathbf a\in R^n$, proves the assertion by induction.  Every tuple
in $(U^{-1}R)^n$ admits a common denominator.

Conversely, clear the finitely many coefficient denominators of a
polynomial over $U^{-1}S$, obtaining $F\in S[\mathbf X]$ whose values
have integral images after localization.  The coefficient classes
generate a finite $R$-submodule $M\subseteq S/R$.  The classes of
the values of $F$ generate a submodule of $M$, hence a finite
submodule.  It vanishes after localization, so one element of $U$
annihilates all its generators.  Multiplication by that element
makes $F$ integer-valued on $R^n$.

For completion, fix a finite rectangular set of monomials, of
cardinality $N$.  The evaluation rows
$(\mathbf a^\alpha)_\alpha$, for $\mathbf a\in R^n$, span a
submodule of $R^N$.  Noetherianity implies that finitely many actual
evaluation rows generate this submodule.  Integer-valuedness within
the chosen degree bound is therefore the kernel of a finite map
\[
 S^N\longrightarrow(S/R)^r.
\]
Tensoring with the flat $R$-algebra $A$ identifies this kernel with
the polynomials over $A\otimes_R S$ having integral values on $R^n$.
If $F$ is such a polynomial, let $M_F$ be the finite $A$-module
generated by its coefficient classes in $(A\otimes_R S)/A$.
For $\mathbf x\in A^n$, choose
$\mathbf a_k\in R^n$ with
$\mathbf x-\mathbf a_k\in(\mathfrak m^kA)^n$.  Then
$[F(\mathbf x)]\in\mathfrak m^kM_F$ for every $k$.
Krull intersection gives $[F(\mathbf x)]=0$.  Taking the union
over the degree bounds proves the completion formula.

Suppose next that $(R,\mathfrak m)$ has finite residue field, and
let $c$ be a coefficient of an integer-valued polynomial.
For each nonmaximal prime $\mathfrak p$, the domain
$R/\mathfrak p$ is infinite.  Rectangular Vandermonde interpolation
on a grid of distinct residues modulo $\mathfrak p$ gives
$v\notin\mathfrak p$ with $vc\in R$; Cramer's rule suffices,
so no coefficient torsion is discarded.  Thus
$\{a\in R:ac\in R\}$ is either $R$ or an ideal with radical
$\mathfrak m$.  It contains a power of $\mathfrak m$, proving
$c\in C$.  If the residue field is infinite, the grid may instead
be chosen with unit differences, giving $c\in R$ directly.

Write $C_j=(R:_S\mathfrak m^j)$.  These submodules contain $R$
and satisfy $C_jC_k\subseteq C_{j+k}$.  Thus $C$ is a subring.
For generators $a_1,\ldots,a_r$ of $\mathfrak m^j$,
\[
 C_j=\ker\bigl(S\longrightarrow(S/R)^r,\ 
                    c\longmapsto(a_ic\bmod R)_i\bigr).
\]
Flat completion preserves this kernel.  Taking the filtered union
proves the asserted formula for $C$.  All identifications arise
from the same polynomial evaluations and multiplications, giving
the stated compatibility.
\end{proof}
\begin{lemma}[Binary injectivity forces a ring epimorphism]
\label{lem:binary-effective-epimorphism}
Let $(A,\mathfrak m)$ be Noetherian local with finite residue field,
let $A\subset C$, and suppose that $C/A$ is killed elementwise by
powers of $\mathfrak m$.  If the binary tensor map for
$\Omega(A,C)$ is injective, then
$A\to C$ is a ring epimorphism.  Equivalently, multiplication
$C\otimes_A C\to C$ is an isomorphism.
\end{lemma}
\begin{proof}
For $c\in C$, choose $j\ge1$ with $\mathfrak m^jc\subseteq A$.
Choose one lift in $A$ of each element of the finite ring
$A/\mathfrak m^j$, and let $P$ be the product of the corresponding
linear factors.  This monic polynomial satisfies
$P(A)\subseteq\mathfrak m^j$, so $P,cP\in\Omega(A,C)$.
The tensor $cP\otimes P-P\otimes cP$ has zero polynomial image.
Injectivity and extraction of the coefficient of $X^{\deg P}$
in each factor give
$c\otimes1=1\otimes c$ in $C\otimes_A C$.
Consequently every $c\otimes d$ equals $cd\otimes1$.
This is the ring-epimorphism criterion
\cite[Lemmas 10.107.1--10.107.3, Tag 04VM]{StacksRingEpimorphisms}.
\end{proof}

\begin{lemma}[The positive local cases]
\label{lem:binary-positive-local}
If $C=A$, every polynomial tensor map is an isomorphism.
The same holds for $C=\operatorname{Frac}(A)$ when $A$ is a
one-dimensional analytically irreducible Noetherian local domain.
\end{lemma}
\begin{proof}
The first case is the ordinary polynomial ring.  In the second case,
Theorem~\ref{thm:int-analytic-free} makes
$M=\operatorname{Int}(A)$ free over $A$.  Let $(f_i)$ be a basis.
After extending scalars to $K=\operatorname{Frac}(A)$ it is a
$K$-basis of $K[X]$, since $A[X]\subseteq M\subseteq K[X]$.
Expand $F\in\operatorname{Int}(A^n)$ in the first variable as
\[
 F(X_1,\ldots,X_n)=\sum_i f_i(X_1)G_i(X_2,\ldots,X_n),
\]
with finitely many terms.  Specializing the other variables to
$A$ gives an element of $M$, so uniqueness of its basis
coefficients shows that each $G_i$ is integer-valued.
Induction gives surjectivity of $\theta_n$.
The products of the $f_i$ in distinct variables are linearly
independent over $K$, and hence over $A$, proving injectivity.
\end{proof}

\subsection{Finite-precision approximation through nilpotents}

For a one-dimensional Cohen--Macaulay local ring, a
\emph{regular parameter} means a nonzerodivisor in the maximal ideal
whose principal ideal is primary to that maximal ideal.
If $A$ is complete and $t$ is such a parameter, then
$Q(A)=A[1/t]$, and the $t$-adic and maximal-ideal topologies agree.

We use the following classical approximation statement:
for a complete one-dimensional Noetherian local domain $D$ with
finite residue field, $\operatorname{Int}(D)$ is uniformly dense in
the continuous functions $D\to D$.  In particular it approximates
every clopen indicator to any prescribed parameter-adic precision.
This is the complete-domain case of
\cite[Proposition 4.3, author version pp.\ 18--19]{CahenChabert1999Skolem};
see also \cite[Section 4, pp.\ 52--53]{CahenGrazziniHaouat1982}.

\begin{lemma}[Indicators with small derivatives]
\label{lem:binary-primary-indicators}
Let $(A,\mathfrak m)$ be a complete one-dimensional Noetherian
Cohen--Macaulay local ring with finite residue field and exactly
one minimal prime $N$.  Put $D=A/N$ and $Q=Q(A)$.
Fix a regular parameter $t$, and write $\bar t$ for its image in $D$.
For every clopen subset $U\subseteq D$ and every $b\ge1$, there is
$H\in Q[X]$ such that, for all $a\in A$,
\[
 H(a)-\mathbf1_U(\bar a)\in t^bA,
 \qquad H'(a)\in t^bA.
\]
In particular $H\in\operatorname{Int}(A)$.
\end{lemma}
\begin{proof}
The ring $Q$ is Artinian local, with nilradical $NQ$ and residue
field $K=\operatorname{Frac}(D)$.  Choose $h\ge1$ with $(NQ)^h=0$;
when $N=0$, take $h=1$.  Choose $e\ge1$ with
$\mathfrak m^e\subseteq tA$.  If $p$ is the residue characteristic,
then
\[
 p^{eM}\in t^MA\qquad(M\ge0),
\]
whether $p$ is regular, nilpotent, or a nonzero zero divisor in $A$.

Approximate $\mathbf1_U$ modulo $\bar tD$ by a polynomial
$f\in K[X]$.  Lift its finitely many coefficients arbitrarily along
$Q\to K$, obtaining $F\in Q[X]$.  For each $a\in A$ write
\[
 F(a)=\chi+tb(a)+n(a),\qquad
 \chi=\mathbf1_U(\bar a)\in\{0,1\},\quad
 b(a)\in A,\quad n(a)\in NQ.
\]
Choose $L,L'\ge0$ such that the coefficients of $F$ lie in
$t^{-L}A$ and $F'(A)\subseteq t^{-L'}A$.  The displayed
decomposition then has $n(a)\in t^{-L}A$ uniformly in $a$.
No ring section $K\to Q$ is used.

Set $q=p^r$.  Choose $r$ so large that all of the following finite
conditions hold:
\begin{align*}
 q&\ge\max\{h,\ b+(L+1)(h-1)\},\\
 r&\ge v_p(j)+e(b+Lj) &&(1\le j<h),\\
 r&\ge v_p(j)+eb &&(1\le j<b),\\
 r&\ge e\bigl(b+L(h-1)+L'\bigr).
\end{align*}
For $1\le j<q$, the integer identity
$j\binom qj=q\binom{q-1}{j-1}$ gives
$v_p(\binom qj)\ge r-v_p(j)$.
If $\chi=0$, nilpotence yields
\[
 F(a)^q=\sum_{0\le j<h}\binom qj
                  (tb(a))^{q-j}n(a)^j,
\]
whose terms lie in $t^{q-j-Lj}A\subseteq t^bA$.
If $\chi=1$, put $v=1+tb(a)$.  Then
\[
 F(a)^q=v^q+\sum_{1\le j<h}\binom qj v^{q-j}n(a)^j.
\]
The second line of bounds makes each nilpotent term lie in $t^bA$.
In the binomial expansion of $v^q-1$, the terms with $j\ge b$
already lie in $t^bA$, while the third line handles the remaining
finitely many terms.  Thus $F(a)^q\equiv\chi\pmod{t^bA}$ uniformly.

For every integer $m\ge0$, the same nilpotent expansion gives
$F(a)^m\in t^{-L(h-1)}A$ uniformly.  Consequently
\[
 (F^q)'(a)=p^rF(a)^{q-1}F'(a)
       \in p^rt^{-L(h-1)-L'}A\subseteq t^bA.
\]
Taking $H=F^q$ proves both conclusions.
\end{proof}

\begin{lemma}[Small reduced values and a unit derivative]
\label{lem:binary-primary-small-values}
Under the hypotheses of Lemma~\ref{lem:binary-primary-indicators},
for every $n\ge1$ there is $g_n\in\operatorname{Int}(A)$ such that
\[
 g_n(A)\subseteq t^nA+N,\qquad
 g_n'(0)\in1+t^nA\subseteq A^\times .
\]
\end{lemma}
\begin{proof}
Choose representatives $d_1,\ldots,d_s$ of the finite ring
$D/\bar t^nD$ and lifts $a_i\in A$.  Apply the preceding lemma
at precision $n$ to the clopen partition
$U_i=d_i+\bar t^nD$, obtaining $H_i$.
Then
\[
 g_n(X)=X-\sum_i a_iH_i(X)
\]
has the required reduced values and derivative.
\end{proof}

\Needspace{8\baselineskip}
\subsection{Nilpotent and multibranch tensor kernels}

\begin{theorem}[Nilpotents and unbounded parameter torsion]
\label{lem:nilpotent-without-section}
Let $(A,\mathfrak m)$ be a complete one-dimensional Noetherian
Cohen--Macaulay local ring with finite residue field.  Put
$M=\operatorname{Int}(A)\subseteq Q(A)[X]$ and let $t$ be any
regular parameter.  If $A$ is nonreduced, then
\[
 \ker\bigl(M\otimes_A M\longrightarrow\operatorname{Int}(A^2)\bigr)
\]
is nonzero.  Every element of this kernel is annihilated by some
power of $t$, but no one power annihilates the whole kernel.
This holds in every characteristic, including when the residue
characteristic is a nonzero zero divisor.
\end{theorem}
\begin{proof}
First suppose that $A$ has a unique minimal prime $N$.
Write $Q=Q(A)$, $D=A/N$, $K=\operatorname{Frac}(D)$, and let
$\rho:Q\to K$ be reduction.  Choose $h\ge2$ with $(NQ)^h=0$.
There is one fixed regular element
\[
 \delta=\prod_{0\le i<j<h}(t^j-t^i)\in A
\]
such that
\begin{equation}\label{eq:binary-uniform-nil-derivative}
 \delta zF'(0)\in A\qquad(z\in N,\ F\in M).                 
\end{equation}
Indeed, $F(zT)$ has degree at most $h-1$.  Its values at
$1,t,\ldots,t^{h-1}$ belong to $A$, and its linear coefficient
is $zF'(0)$.  The adjugate of this fixed Vandermonde matrix gives
\eqref{eq:binary-uniform-nil-derivative}.  Every factor of $\delta$ is a power of $t$ times a unit,
so $\delta$ is regular.  In particular,
\begin{equation}\label{eq:binary-nil-submodule-derivative}
 \delta G'(0)\in A\qquad(G\in NM).                         
\end{equation}

Choose a nonzero socle element of the Artinian local ring $Q$.
Its socle lies in $NQ$, since $Q$ is nonreduced.
After multiplication by a power of $t$, it gives $u\in A$ with
\[
 u\ne0,\qquad uNQ=0,\qquad u^2=0,\qquad
 \operatorname{Ann}_A(u)=N.
\]
For $\lambda\in K$, multiplication of $u$ by any lift of
$\lambda$ to $Q$ is independent of the lift.  Thus
\[
 L_u=\{\lambda\in K:u\lambda\in A\}
\]
is a $D$-module.  Multiplication by $u$ identifies it with
$A\cap uQ$, a finite $A$-module annihilated by $N$.
Hence $L_u$ is a finite $D$-submodule of $K$.
Choose $0\ne b\in D$ with $bL_u\subseteq D$, and choose a lift
$\widetilde b\in A$.  This lift is regular.

For $F\in M$, evaluation at the square-zero element $u$ gives
$uF'(0)=F(u)-F(0)\in A$, so $\rho(F'(0))\in L_u$.
Therefore
\begin{equation}\label{eq:binary-nil-detector}
 \Phi:M\otimes_A M\longrightarrow M/NM,\qquad
 \Phi(F\otimes G)=b\rho(F'(0))[G]                          
\end{equation}
is well-defined.  The target is a $D$-module, and the scalar in
\eqref{eq:binary-nil-detector} belongs to $D$.  Balancing is explicit: for $a\in A$,
\[
 b\rho((aF)'(0))[G]
 =b\rho(a)\rho(F'(0))[G]
 =b\rho(F'(0))[aG].
\]
Thus the detector is balanced over $A$ itself.

Take $g_n$ from Lemma~\ref{lem:binary-primary-small-values} and put
\[
 h_n=\frac{u}{t^n}g_n,\qquad
 w_n=h_n\otimes g_n-g_n\otimes h_n.
\]
The reduced-value condition gives $g_n(A)\subseteq t^nA+N$.
Since $uN=0$, one has $h_n(A)\subseteq uA\subseteq A$.
Both tensor factors therefore lie in $M$.  Their antisymmetric
tensor has zero polynomial image because $h_n$ is a scalar
multiple of $g_n$ in $Q[X]$.

Fix $c\ge0$.  If $t^cw_n=0$, apply \eqref{eq:binary-nil-detector}.  Since
$\rho(h_n'(0))=0$, this gives
\[
 \widetilde b\,t^c g_n'(0)h_n\in NM.
\]
Differentiation at zero and \eqref{eq:binary-nil-submodule-derivative} yield
\[
 \delta\widetilde b\,u\,t^{c-n}(g_n'(0))^2\in A.
\]
The derivative $g_n'(0)$ is a unit of $A$, so, for $n>c$,
\begin{equation}\label{eq:binary-nil-krull}
 \delta\widetilde b\,u\in t^{n-c}A.                       
\end{equation}
The left side is a fixed nonzero element, because $\delta$ and
$\widetilde b$ are regular.  Krull intersection shows that \eqref{eq:binary-nil-krull}
fails for all sufficiently large $n$.  Thus $t^cw_n\ne0$ for
all sufficiently large $n$.

Now allow several minimal primes.  Since $A$ is one-dimensional
Cohen--Macaulay, its associated primes are its minimal primes.
Choose the isolated primary components $\mathfrak q_i$ of zero
and put $A_i=A/\mathfrak q_i$.  Then
\[
 A\hookrightarrow\prod_iA_i,\qquad
 Q(A)=\prod_iQ(A_i).
\]
Each $A_i$ is complete, primary, one-dimensional Cohen--Macaulay,
and has the same finite residue field.  The image $t_i$ of $t$
is a regular parameter.  The cokernel of the displayed inclusion
is finite and supported only at $\mathfrak m$: localization at
each minimal prime leaves only the corresponding factor and
gives an isomorphism.  Consequently, for one fixed $d\ge0$,
\begin{equation}\label{eq:binary-primary-conductor}
 t^d\prod_iA_i\subseteq A.                               
\end{equation}
If $A$ is nonreduced, at least one $A_i$ is nonreduced.
Let $e_i$ be its generic idempotent and let $w_n$ be the tensors
constructed for $A_i$ above.  The polynomials
\[
 \widetilde h_n=t^de_i h_n,\qquad
 \widetilde g_n=t^de_i g_n
\]
preserve $A$ by \eqref{eq:binary-primary-conductor}.  Their antisymmetric tensor has zero image.
Projection of coefficients gives a map
$\operatorname{Int}(A)\to\operatorname{Int}(A_i)$, since every
input in $A_i$ lifts to $A$.  The induced map of tensor products
sends the supported tensor to $t_i^{2d}w_n$.  More generally its
$t^c$-multiple maps to $t_i^{c+2d}w_n$, which is nonzero for
large $n$.  This proves the required unboundedness for $A$.

Finally $M[1/t]=Q(A)[X]$, so the localized polynomial tensor map
is an isomorphism.  Every individual kernel element is therefore
annihilated by some power of $t$, as asserted.
\end{proof}

\begin{lemma}[Transfer through a regular conductor]
\label{lem:binary-selected-factor-transfer}
Let $A\subset C$ be a commutative inclusion, and let $e\in C$
be an idempotent such that $eC=Q(A_0)$, where $A_0=eA$ is a
complete one-dimensional Cohen--Macaulay local ring with finite
residue field.  Suppose that $A_0$ is nonreduced and, for a
regular parameter $t_0$ of $A_0$ and some $d\ge0$,
\begin{equation}\label{eq:binary-selected-conductor}
 t_0^deA_0\subseteq A.                                   
\end{equation}
Then the binary tensor map for $\Omega(A,C)$ is not injective.
The same conclusion holds if \eqref{eq:binary-selected-conductor} is obtained from a finite full
order $A_0\subseteq O\subseteq Q(A_0)$ and a regular conductor
$a eO\subseteq A$, with $a$ a nonzerodivisor of $O$.
If this conductor lies in an ideal of $A$, the supported
polynomials in the proof have all their values in that ideal.
\end{lemma}
\begin{proof}
Projection gives an $A$-linear map
$\Omega(A,C)\to\operatorname{Int}(A_0)$ and hence a map from its
binary tensor product to
$\operatorname{Int}(A_0)\otimes_{A_0}\operatorname{Int}(A_0)$.
By Theorem~\ref{lem:nilpotent-without-section}, choose an
antisymmetric kernel tensor
$w=h\otimes g-g\otimes h$ whose $t_0^{2d}$-multiple is nonzero.
The supported polynomials $t_0^deh$ and $t_0^deg$ preserve $A$
by \eqref{eq:binary-selected-conductor}.  Their antisymmetric tensor has zero polynomial image,
while projection sends it to $t_0^{2d}w\ne0$.

For the order formulation, $O/aO$ has finite length as an
$A_0$-module.  Thus $t_0^dO\subseteq aO$ for some $d$,
which implies \eqref{eq:binary-selected-conductor} and preserves any specified ideal containing
the conductor.  No tensor-preservation of injections is used.
\end{proof}

\begin{lemma}[Several reduced branches give a binary kernel]
\label{lem:binary-reduced-branches}
Let $A$ be a complete reduced one-dimensional Noetherian local
ring with finite residue field.  If $A$ has more than one minimal
prime, the binary tensor map for $\operatorname{Int}(A)$ is not
injective.
\end{lemma}
\begin{proof}
Such a ring is Cohen--Macaulay.  Put $D_i=A/\mathfrak p_i$
for its minimal primes, and identify $A$ with its image in
$\prod_iD_i$.  Let $e_i$ be the corresponding generic idempotents.
The conductor argument \eqref{eq:binary-primary-conductor} shows that
\[
 I_i=\{a\in D_i:e_i a\in A\}
\]
is a nonzero ideal of $D_i$.  Choose $0\ne v_i\in I_i$ and
$u_i\in I_i\setminus\mathfrak m_iI_i$.
Polynomial approximation on the complete domain $D_i$ gives
$f_i\in\operatorname{Int}(D_i)$ such that
\[
 f_i(v_i)-f_i(0)\equiv1\pmod{\mathfrak m_i}.
\]
Indeed a sufficiently small clopen neighborhood of $v_i$ misses
zero, and its indicator can be approximated modulo a parameter.
Then $H_i=e_i u_i f_i$ belongs to $\operatorname{Int}(A)$.

For $F\in\operatorname{Int}(A)$ write $F_i$ for its $i$th
component.  The difference
\[
 \Delta_i(F)=F_i(v_i)-F_i(0)
\]
belongs to $I_i$, because $F(e_i v_i)-F(0)\in A$ is supported
only on component $i$.  This defines an $A$-linear map to $I_i$.
Compose with a residue-field linear functional
$I_i/\mathfrak m_iI_i\to\kappa$, where $\kappa=A/\mathfrak m$,
taking $[u_i]$ to $1$.  The resulting map
$r_i:\operatorname{Int}(A)\to\kappa$ satisfies $r_i(H_i)=1$.
For distinct $i,j$, the tensor $H_i\otimes H_j$ has zero
polynomial image, but $r_i\otimes r_j$ sends it to
$1\in\kappa\otimes_A\kappa=\kappa$.  It is therefore nonzero.
\end{proof}

\begin{corollary}[Complete local curves]
\label{cor:binary-complete-cm}
For a complete one-dimensional Noetherian Cohen--Macaulay local
ring $A$ with finite residue field, the binary tensor map for
$\operatorname{Int}(A)$ is injective if and only if $A$ is a
domain.  In that case all its polynomial tensor maps are
isomorphisms.
\end{corollary}
\begin{proof}
The nonreduced and reduced multibranch cases are excluded by the
two preceding obstructions.  A complete local domain is
analytically irreducible, so the converse follows from
Lemma~\ref{lem:binary-positive-local}.
\end{proof}

\subsection{The complete local coefficient ring}

\begin{lemma}[A complete local epimorphic extension]
\label{lem:complete-epic-decomposition}
Let $(A,\mathfrak m)$ be complete Noetherian local with finite
residue field.  Let $A\subset C$ be a ring epimorphism such that
$C/A$ is killed elementwise by powers of $\mathfrak m$.  Define
\[
 D=\bigcap_{j\ge0}\mathfrak m^jC,\qquad I=A\cap D.
\]
Then
\[
 C=A+D,\qquad D=\mathfrak mD=D^2=eC
\]
for an idempotent $e\in C$.  In particular $D$ is a finitely
generated $C$-ideal, and
\[
 C\simeq D\times A/I,\qquad a\longmapsto(ea,a+I)
\]
is an isomorphism of $A$-algebras.
\end{lemma}
\begin{proof}
A nonzero epimorphic algebra over a field is that field.
Indeed, $b\otimes1=1\otimes b$, and applying a linear functional
which takes $1$ to $1$ in one factor forces $b$ to be a scalar.
Since ring epimorphisms are preserved by base change, either
$C/\mathfrak mC=0$ or $C/\mathfrak mC=A/\mathfrak m$.
In the first case $\mathfrak mC=C$, so $D=C$, $I=A$, and $e=1$.

In the other case $\mathfrak n=\mathfrak mC$ is maximal and
$A\to C_{\mathfrak n}$ is a local ring epimorphism.
Ferrand's complete-local epimorphism lemma
\cite[Lemma 2.1, p.\ 5]{Ferrand1968NoetherianMonomorphisms}
implies that it is surjective.  The lemma requires the source to
be complete Noetherian local; it imposes no Noetherian hypothesis
on the target.  Write $C_{\mathfrak n}=A/J$, with $J\subseteq
\mathfrak m$.

The epimorphism $A/J\to C/JC$ has the retraction induced by
localization to $C_{\mathfrak n}$.  A split monomorphism which is
also a ring epimorphism is an isomorphism: if $i$ and $r$ are the
section and retraction, then $ir$ and the identity agree after
precomposition with the epimorphism $i$.
Consequently
\[
 C/JC=A/J,\qquad C=A+JC,\qquad A\cap JC=J.
\]
The ideal $J$ has finitely many $A$-generators, all zero in
$C_{\mathfrak n}$.  Multiplying finitely many localization
denominators gives $v\in C\setminus\mathfrak n$ with $vJC=0$.
Its image in $C/JC=A/J$ is a unit.  Choose $w\in C$ so that
$e=1-vw\in JC$.  For $z\in JC$, one has $z=ez$, whence
$e^2=e$ and $JC=eC$.  Thus
\[
 C=eC\times(1-e)C,\qquad (1-e)C=A/J.
\]
Because $e\in JC$ and $J\subseteq\mathfrak m$, multiplication
by $e$ in an expression for $e$ shows
$\mathfrak m(eC)=eC$.  Krull intersection on the finite
$A$-module $A/J$ now gives
\[
 \bigcap_{j\ge0}\mathfrak m^jC
   =eC\times\bigcap_{j\ge0}\mathfrak m^j(A/J)
   =eC\times0.
\]
Therefore $D=eC$ and $I=J$, proving every assertion, including
the named intersection and the finite generation of $D$.
\end{proof}

\begin{lemma}[An empty closed fiber is a punctured curve]
\label{lem:binary-empty-closed-fiber}
Let $(B,\mathfrak m_B)$ be a nonzero Noetherian local ring and
let $B\subset T$ satisfy
$T/B$ elementwise $\mathfrak m_B$-power torsion and
$\mathfrak m_BT=T$.  Then $B$ is one-dimensional
Cohen--Macaulay and $T=Q(B)$ canonically as a $B$-algebra.
\end{lemma}
\begin{proof}
For every $x\in\mathfrak m_B$, localization gives $B_x=T_x$.
If $x_1,\ldots,x_r$ generate $\mathfrak m_B$, their principal
opens cover $\operatorname{Spec}(T)$, because
$(x_1,\ldots,x_r)T=T$.  These canonical identifications glue to
\[
 \operatorname{Spec}(T)
 \simeq\operatorname{Spec}(B)\setminus\{\mathfrak m_B\}.
\]
The punctured spectrum is affine.  The affine-puncture criterion
\cite[Lemma 31.17.1, Tag 0BCR]{StacksAffinePuncture}
gives $\dim B\le1$.  Dimension zero is impossible, since then
$\mathfrak m_B$ is nilpotent while $\mathfrak m_BT=T\ne0$.

If $\mathfrak m_B^j b=0$, use
$1\in\mathfrak m_B^jT$ to obtain $b=0$ in $T$, and hence in $B$.
Thus the maximal ideal is not an associated prime of $B$.
It follows that $B$ has depth one and is Cohen--Macaulay.
For a regular parameter $t$, the punctured spectrum equals
$D(t)$, so the canonical isomorphism above gives
$T=B[1/t]=Q(B)$.  The last equality follows because every
regular nonunit has a maximal-ideal-primary principal ideal.
\end{proof}

\begin{lemma}[The selected component and its conductor ideal]
\label{lem:selected-punctured-curve}
In Lemma~\ref{lem:complete-epic-decomposition}, suppose that
$D\ne0$ and put $A_0=eA\subset D$.
Then $A_0$ is a complete one-dimensional Cohen--Macaulay local
ring and $D=Q(A_0)$.  The ideal $I$ embeds as an ideal of $A_0$
containing a nonzerodivisor.  If $I$ is proper in $A$, its image
is proper in $A_0$.
\end{lemma}
\begin{proof}
The image $A_0$ is a complete Noetherian local quotient of $A$.
The quotient $D/A_0$ is maximal-ideal-power torsion and
$\mathfrak m_0D=D$.  Lemma~\ref{lem:binary-empty-closed-fiber}
therefore gives the dimension, depth, and total quotient ring.
For $i\in I$, the identity $ei=i$ proves injectivity of
$I\to A_0$ and identifies its image as an ideal.
Since $D/I=C/A$, the class of $e$ is killed by a power of
$\mathfrak m$.  A corresponding power of any regular parameter
of $A_0$ therefore lies in the image of $I$.
If this image were all of $A_0$, then $e\in I\subseteq A$.
A nonzero idempotent in the local ring $A$ equals $1$,
which would give $I=A$.
\end{proof}

\begin{lemma}[A proper closed-fiber factor gives a binary kernel]
\label{lem:binary-proper-closed-fiber}
Under the hypotheses of Lemma~\ref{lem:complete-epic-decomposition},
if $D\ne0$ and $I\ne A$, the binary tensor map for $\Omega(A,C)$
is not injective.
\end{lemma}
\begin{proof}
Put $B=eA$, $T=A/I$, and $J=\ker(A\to B)$.  Then
$I\cap J=0$.  The image $L$ of $I$ in $B$ is proper and
contains a regular element by
Lemma~\ref{lem:selected-punctured-curve}.  Thus $L$ is primary
to the maximal ideal of the one-dimensional ring $B$, and
\[
 E=B/L=A/(I+J)
\]
is a nonzero finite Artinian local ring.  The usual two-ideal
fiber-product identity gives the original inclusion in the form
\begin{equation}\label{eq:binary-closed-fiber-product}
 A=B\times_E T\ \subseteq\ Q(B)\times T=C.                 
\end{equation}
All these local rings have residue field $\kappa=A/\mathfrak m$.

Let $B_i$ be the isolated primary factors of $B$, with generic
idempotents $\varepsilon_i$ in $Q(B)$, and put
\[
 L_i=\{z\in B_i:\varepsilon_i z\in L\}.
\]
Choose a regular parameter $a\in L$.  The finite
primary-component conductor gives
$a^d\prod_iB_i\subseteq B$ for some $d$; multiplying by $a$
shows that
$a_i^{d+1}B_i\subseteq L_i$.  Hence each $L_i$ is a nonzero
finite ideal.  Fix one $i$, put $v_i=a_i^{d+1}$, and choose
$u_i\in L_i\setminus\mathfrak m_iL_i$.
The reduced image of $v_i$ is nonzero.  A clopen neighborhood
separating it from zero, together with
Lemma~\ref{lem:binary-primary-indicators}, gives
$f_i\in\operatorname{Int}(B_i)$ with
\[
 f_i(v_i)-f_i(0)\equiv1\pmod{\mathfrak m_i}.
\]
The polynomial
$H=(\varepsilon_i u_i f_i,0)$ preserves $A$, since its first
component has values in $L$ and its second component is zero.

The element $z=(\varepsilon_i v_i,0)$ belongs to $A$.
For $F\in\Omega(A,C)$, the difference $F(z)-F(0)$ lies in $A$
and is supported only on component $i$ of the first factor
in \eqref{eq:binary-closed-fiber-product}.  Therefore
$\Delta_i(F)=F_i(v_i)-F_i(0)\in L_i$.
This is $A$-linear.  Composing with a $\kappa$-linear functional
$L_i/\mathfrak m_iL_i\to\kappa$ taking $[u_i]$ to $1$ gives
\[
 r:\Omega(A,C)\longrightarrow\kappa,\qquad r(H)=1.
\]

Choose in $T$ one lift $\widetilde e$ of each $e\in E$, and set
\[
 P(X)=\prod_{e\in E}(X-\widetilde e)\in T[X].
\]
It is monic of degree $|E|$ and has values in $\ker(T\to E)$.
Thus $K=(0,P)\in\Omega(A,C)$.
Taking the coefficient of $X^{|E|}$ in the $T$-component and
reducing to $\kappa$ gives an $A$-linear map $s$ with $s(K)=1$.
The tensor $H\otimes K$ has zero polynomial image because
its coefficient components are orthogonal.  However,
$r\otimes s$ sends it to
$1\in\kappa\otimes_A\kappa=\kappa$, so it is nonzero.
This construction also applies when the selected primary factor
$B_i$ is nonreduced.
\end{proof}

\begin{proposition}[Complete local classification]
\label{prop:complete-arbitrary-coefficients}
Let $(A,\mathfrak m)$ be complete Noetherian local with finite
residue field, and let $A\subset C$ have $C/A$ killed elementwise
by powers of $\mathfrak m$.  The binary tensor map for
$\Omega(A,C)$ is injective if and only if
\[
 C=A\quad\text{or}\quad
 \bigl(A\text{ is a one-dimensional domain and }
                   C=\operatorname{Frac}(A)\bigr).
\]
In either case all polynomial tensor maps are isomorphisms and
the coaddition and comultiplication preimages are unique.
\end{proposition}
\begin{proof}
Binary injectivity makes $A\to C$ a ring epimorphism by
Lemma~\ref{lem:binary-effective-epimorphism}.
Use the notation of Lemma~\ref{lem:complete-epic-decomposition}.
If $D=0$, then $C=A$.
If $D\ne0$, Lemma~\ref{lem:binary-proper-closed-fiber}
rules out $I\ne A$.  Hence $I=A$, $e=1$ and $C=D$.
Lemma~\ref{lem:selected-punctured-curve} now gives
$\dim A=\operatorname{depth}A=1$ and $C=Q(A)$.
Corollary~\ref{cor:binary-complete-cm} forces $A$ to be a domain.
The converse and the all-arity assertion follow from
Lemma~\ref{lem:binary-positive-local}.
\end{proof}

\begin{corollary}[A completed domain with arbitrary coefficient torsion]
\label{cor:binary-completed-domain}
Let $(R,\mathfrak m_R)$ be a nonfield Noetherian local domain
with finite residue field, let $K=\operatorname{Frac}(R)$, and
put $A=\widehat R$ and $Q=K\otimes_R A$.
If $A\subset C$ and $C/A$ is elementwise killed by powers of
the maximal ideal of $A$, then binary injectivity, all-arity
tensor isomorphisms, and unique polynomial operations are
equivalent.  They hold exactly when
\[
 C=A\quad\text{or}\quad
 \bigl(\dim R=1,\ A\text{ is a domain},\ C=Q\bigr),
\]
with the second equality understood as an equality of
canonical $A$-algebras.
\end{corollary}
\begin{proof}
Apply Proposition~\ref{prop:complete-arbitrary-coefficients}.
In its nontrivial case $\dim R=\dim A=1$.
Choose $0\ne a\in\mathfrak m_R$.
Its image is nonzero in the one-dimensional local domain $A$,
so $A[1/a]=\operatorname{Frac}(A)$.
Therefore $K\otimes_R A=\operatorname{Frac}(A)$ canonically.
This proves the exact stated form.  The other implications
follow from the proposition and uniqueness for the zero
polynomial.
\end{proof}

\subsection{Canonical descent and globalization}

\begin{proof}[Proof of Theorem~\ref{thm:binary-injective-classification}]
Suppose first that $\theta_2$ is injective.
Lemma~\ref{lem:arbitrary-coefficient-basechange} localizes this
condition at every maximal ideal.  At an infinite-residue
maximal ideal the integer-valued ring is the ordinary polynomial
ring, so there is no further restriction.

Fix a finite-residue maximal ideal, write
$A=R_{\mathfrak m}$ and $S'=S_{\mathfrak m}$, and let $C$ be
the stated effective coefficient subring of $S'$.
The coefficient assertion of the base-change lemma gives
$\Omega_n(A,S')=\Omega_n(A,C)$ for every $n$.
Set $\widehat C=\widehat A\otimes_A C$ and
$\widehat S'=\widehat A\otimes_A S'$.
Flatness embeds $\widehat A\subseteq\widehat C\subseteq\widehat S'$,
and the same lemma identifies the completed polynomial tensor
maps with those over $\widehat A$.
The quotient $\widehat C/\widehat A$ is elementwise killed by
powers of the maximal ideal.  Proposition
\ref{prop:complete-arbitrary-coefficients} applies.

If $\widehat C=\widehat A$, faithful flatness gives $C=A$.
Otherwise $\widehat A$ is a one-dimensional domain and
$\widehat C=\operatorname{Frac}(\widehat A)$.
The original $A$ embeds in its completion, so it is a domain;
dimension invariance under completion gives $\dim A=1$.
It remains to descend the coefficient subring.
For every $0\ne a\in A$, multiplication by $a$ on
$\widehat C$ is bijective.  Its kernel and cokernel are the
faithfully flat base changes of those of multiplication by
$a$ on $C$, so $a:C\to C$ is bijective.
In particular there is an element $c_a\in C$ with $ac_a=1$.
Thus all nonzero elements of $A$ have their unique inverses
inside $C\subseteq S'$, giving the canonical embedding
\[
 \operatorname{Frac}(A)\subseteq C\subseteq S'.
\]
Conversely, for $c\in C$, choose $j$ with
$\mathfrak m_A^jc\subseteq A$ and a nonzero
$a\in\mathfrak m_A^j$; for $j=0$ use $a=1$.
Then $c=(ac)/a$ lies in that canonical fraction field.
Hence $C=\operatorname{Frac}(A)$, proving condition
\textup{(iv)} without any initial torsion-free hypothesis on $S'$.

Conversely, assume \textup{(iv)}.
At every finite-residue maximal ideal, the localized polynomial
maps are isomorphisms by Lemma~\ref{lem:binary-positive-local};
at every infinite-residue maximal ideal they are ordinary
polynomial tensor maps.  Localization detects the kernels and
cokernels of arbitrary module homomorphisms.  Thus every global
$\theta_n$ is an isomorphism, proving \textup{(ii)}.
Both substitutions $F(X+Y)$ and $F(XY)$ belong to
$\Omega_2(R,S)$, so \textup{(ii)} implies \textup{(iii)}.
As observed above, uniqueness for $F=0$ gives
\textup{(iii)}$\Rightarrow$\textup{(i)}.

The tensor isomorphisms transport polynomial coaddition,
comultiplication, scalar operations and the evaluations at
$0$ and $1$ to $\Omega(R,S)$.
Their identities hold after applying the appropriate $\theta_n$,
and injectivity gives the same identities before applying it.
Composition also preserves $\Omega(R,S)$, because both
polynomials take $R$ into $R$.
Transporting these composition identities gives the compatible
plethory structure.  Every compatible operation has the stated
polynomial image, and the tensor maps are injective, so this
structure is unique.
\end{proof}

\begin{proof}[Proof of Theorem~\ref{thm:arbitrary-noetherian-coefficients}]
The assertion is the full four-way equivalence in
Theorem~\ref{thm:binary-injective-classification}, together with
its plethory conclusion.
\end{proof}

\begin{corollary}[Domain bases with coefficient torsion]
\label{cor:domain-torsion-coefficients}
Let $R$ be a nonfield Noetherian domain with fraction field $K$,
and let $R\subset S$ be any unital commutative inclusion.
Binary injectivity, all-arity tensor isomorphisms, and unique
polynomial operations are equivalent.  They hold exactly when
every finite-residue maximal ideal $\mathfrak m$ satisfies
\[
 C_{\mathfrak m}=R_{\mathfrak m}
 \quad\text{or}\quad
 \bigl(\dim R_{\mathfrak m}=1,\
       \widehat{R_{\mathfrak m}}\text{ is a domain},\
       C_{\mathfrak m}=K\bigr).
\]
The equality with $K$ specifies its canonical copy in
$S_{\mathfrak m}$.  No torsion-free hypothesis on $S$ is needed.
\end{corollary}
\begin{proof}
The fraction field of any localization of a domain is $K$,
so this is precisely the domain specialization of
Theorem~\ref{thm:binary-injective-classification}.
\end{proof}

\begin{corollary}[A selected factor with a nonzero conductor]
\label{cor:binary-selected-nonzero-conductor}
Let $(A,\mathfrak m)$ be complete Noetherian local with finite
residue field, and let $A\subset C$ be a unital commutative
inclusion.  Suppose that an idempotent $e\in C$ selects a
nonreduced local Artinian ring $Q_0=eC$, that $A_0=eA$ is
one-dimensional Cohen--Macaulay, and that a finite full order
\[
 A_0\subseteq O\subseteq Q_0,\qquad Q(O)=Q_0,
\]
has a nonzero conductor into $A$:
\[
 \mathfrak c=\{x\in O:e xO\subseteq A\}\ne0.
\]
Then the binary tensor map for $\Omega(A,C)$ is not injective.
Here the conductor is only required to be nonzero; it may
consist entirely of zero divisors.
\end{corollary}
\begin{proof}
The image $A_0$ is a complete local quotient of $A$ with the
same finite residue field.  A regular parameter $t_0$ of $A_0$
is not nilpotent.  Its image in the Artinian local ring $Q_0$
is therefore a unit.  Every regular element of $A_0$ is likewise
a unit in $Q_0$, so $Q(A_0)$ embeds in $Q_0$.
The Artinian ring $Q(A_0)$ cannot have a nontrivial idempotent,
because $Q_0$ has none.  Hence $A_0$ has exactly one minimal
prime.

The finite $A_0$-algebra $O$ is $t_0$-torsion-free, and
$O[1/t_0]$ is Artinian.  Every regular element of $O$ remains
regular there and hence becomes a unit.  Thus
$O[1/t_0]=Q(O)=Q_0$.
The conductor $\mathfrak c$ is an ideal of $O$, is contained
in $A_0$, and remains nonzero after inverting $t_0$.
A nonzero ideal of a local Artinian ring meets its socle.
Choose such a socle element in $\mathfrak c[1/t_0]$ and clear
a power of $t_0$, obtaining
\[
 0\ne u\in\mathfrak c\cap\operatorname{Soc}(Q_0).
\]
The socle of the nonreduced $Q_0$ is contained in its nilradical,
so $u^2=0$.  It also annihilates the nilradical of the subring
$Q(A_0)$.

Apply Lemma~\ref{lem:binary-primary-small-values} to $A_0$ and
$t_0$, and define
\[
 h_n=\frac{u}{t_0^n}g_n\in Q(A_0)[X].
\]
As before, $h_n(A_0)\subseteq uA_0\subseteq\mathfrak c$.
The supported polynomial $e h_n$ therefore belongs to
$\Omega(A,C)$.  For arbitrarily large $n$, it does not belong
to $A[X]$: otherwise projection of its derivative at zero
would give
$u g_n'(0)/t_0^n\in A_0$, hence
$u\in t_0^nA_0$, contradicting separatedness since $g_n'(0)$
is a unit.

Thus $\Omega(A,C)\ne A[X]$, while $eu\in A$ is a nonzero
square-zero element.  If the binary tensor map were injective,
Theorem~\ref{thm:binary-injective-classification}, applied to
this complete local base, would force either the ordinary
polynomial case or a domain base.  Both conclusions contradict
what was just proved.
\end{proof}
